% 第 7 题三视图：按固定源内容流矢量层精确值复刻（1 pt → 0.03515 cm，与全库同制）。
% 原点取正视图左下顶点（源页 (517.1250, 198.8984) pt），y 向上；三幅足尺 62.3398 pt = 1 单位。
% 源实测：正视图＝直角三角形（直角在右下）、侧视图＝同向直角三角形（直角亦在右下，非镜像）、
% 俯视图＝正方形＋自右下至左上的一条实线对角线；三组「双向箭头+端刻度线」尺寸标注（数字 1）；
% 全图无虚线、无填充（除 6 个尺寸箭头）。题注括号按全库先例用全角。
\begin{tikzpicture}[x=0.03515cm,y=0.03515cm,>=stealth,every node/.style={font=\small,inner sep=0pt},
  gl/.style={line width=0.398bp,line cap=round,line join=round}]
% ===== 三幅轮廓 =====
\draw[gl] (0,0)--(62.3398,0)--(62.3398,62.3398)--cycle;            % 正（主）视图（直角右下）
\draw[gl] (93.5117,0)--(155.8516,0)--(155.8516,62.3398)--cycle;     % 侧（左）视图（同向）
\draw[gl] (0,-93.5078)--(62.3398,-93.5078)--(62.3398,-31.1680)--(0,-31.1680)--cycle; % 俯视图
\draw[gl] (62.3398,-93.5078)--(0,-31.1680);                        % 俯视图对角线
% ===== 尺寸标注（a）正视图宽度 y=-6.2344 =====
\draw[gl] (0,-3.1172)--(0,-9.3516);
\draw[gl] (62.3398,-3.1172)--(62.3398,-9.3516);
\draw[gl,->] (24.9375,-6.2344)--(0,-6.2344);
\draw[gl,->] (37.4063,-6.2344)--(62.3398,-6.2344);
\node at (31.3800,-6.9666) {1};
% =====（b）侧视图宽度 =====
\draw[gl] (93.5117,-3.1172)--(93.5117,-9.3516);
\draw[gl] (155.8516,-3.1172)--(155.8516,-9.3516);
\draw[gl,->] (118.4453,-6.2344)--(93.5117,-6.2344);
\draw[gl,->] (130.9141,-6.2344)--(155.8516,-6.2344);
\node at (124.9200,-6.9666) {1};
% =====（c）竖向高度 x=68.5742（数字落在轴线断口）=====
\draw[gl] (65.4570,62.3398)--(71.6914,62.3398);
\draw[gl] (65.4570,0)--(71.6914,0);
\draw[gl,->] (68.5742,40.5234)--(68.5742,62.3398);
\draw[gl,->] (68.5742,21.8203)--(68.5742,0);
\node at (68.8000,30.4534) {1};
% ===== 三幅题注 =====
\node at (31.3750,-19.4266) {正（主）视图};
\node at (124.9250,-19.4266) {侧（左）视图};
\node at (31.3750,-100.6516) {俯视图};
\end{tikzpicture}
